\documentclass[fleqn,a5paper,12pt,twoside]{guitartabs}
\geometry{left=20mm, top=10mm, right=10mm, bottom=10mm,nohead,nofoot}

%Russian-specific packages
%--------------------------------------
\usepackage[T2A]{fontenc}
\usepackage[utf8]{inputenc}
\usepackage[russian]{babel}
\usepackage{sectsty}
%--------------------------------------
\usepackage{subcaption}[]%для выравнивания картинок справа и слева
\usepackage[rightcaption]{sidecap}%тоже для картиночек
\usepackage{setspace}%для ужимки текста

%---------------------------------------

%-----------------------------
%в пакете gtrcrd вместо H используется B
%От лица себя скажу что это не удобно, поэтому я скачал пакет и заменил все B на H
\usepackage{gtrcrd} % Words & Chords edition.
\usepackage{guitarchordschemes} %аппликатуры аккордов
\newcommand{\musSharp}{\texttt{\#}}
\newcommand{\musVspace}{\vspace{2mm}}
\newenvironment{musRepeat}{\left.\begin{matrix*}[l]}
{\end{matrix*}\right\}\textup{2 раза}}
\usepackage{mathtools}  %костыли для скобок
\let\originalleft\left
\let\originalright\right
\renewcommand{\left}{\mathopen{}\mathclose\bgroup\originalleft}
\renewcommand{\right}{\aftergroup\egroup\originalright}%костыли по поводу пробелов вокруг скобок
\newrobustcmd*\muschordscheme[1][]{\chordscheme[#1]\hspace{6mm}} %костыли по выравниванию аккордов
\begin{document}
\setlength{\mathindent}{0pt}%костыли
\setlength{\jot}{0mm}%костыли для межстрочного интервала в математическом режиме
\setlength{\parindent}{0mm}%убираем абзацные отступы
\captionsetup[subfigure]{labelformat=empty} %убираем надписи под картинками
\sectionfont{\LARGE}
\sffamily

\setcounter{page}{0} % начать нумерацию с номера 0
Версия песенника 3.1
\begin{enumerate}
\item Добавлены и удалены все песни по правкам от 28.03.2023 16:55 Кроме:
\begin{itemize}
\item «Не бывает» – жаворонки (на согласовании)
\item Земной простор - походные (опечатка, есть в следопытских)
\end{itemize}
\item оглавление расширено ещё на одну страницу, т.к. стало слишком большим. (возможно вернуть размер перестановкой разделов и рассечением одного раздела на развороте)
\item некоторые песни сов и лирики переставлены местами, чтобы выше были более популярные, а так же чтобы на развороте находились более сочетаемые песни.
\item версия включает исправления редакции от 25.12.2022. Есть спорные моменты, см. ответный документ верстальщика
\end{enumerate}
Необходимо далее:
\begin{enumerate}
\item Проверить грамотность. (в том числе активизировать воинов буквы ё. Верстальщик за этим не следил)
\item Уделить внимание самим аккордам. Не факт что все песни в правильных/удобных тональнастях, нужно призвать для этого гитаристов
\item ещё раз проверить и согласовать текст последней страницы
\item ровнять отступы, куплеты с припевами, переставлять местами песни ради лучшей читаемость (работа верстальщика)
\end{enumerate}
\newpage
Версия песенника 3.2
\begin{enumerate}
\item Добавлены и удалены все песни по правкам от 28.03.2023 и
дополнительно Туман - Максим Леонидов
\item оглавление расширено на одну страницу, разделы сов и жаворонков поменялись местами, чтобы лучше расположиться в оглавлении. Оглавление Жаворонков рассечено на развороте (как же мне не нравится это решение)
\item некоторые песни всех разделов переставлены местами, чтобы на развороте были более визуально сочетаемые песни.
Указание песни в справочных материалах так же изменено
в соответсвии с новыми страницами
\item версия включает исправления редакции от 25.12.2022. Есть
спорные моменты, см. ответный документ верстальщика
(внимание! в новой версии песни в другом порядке. Страницы из прошлой редакии не актуальны, извините за неудобства)
\item раздел “теория” был уточнён и исправлен, консультируясь
с 2 людьми с муз. образованием
\item исправлена финальная страница с благодарностями. Добавилась страница для заметок (вынужденно из-за чётности)
\end{enumerate}
Необходимо далее:
\begin{enumerate}
\item Проверить грамотность. (в том числе активизировать воинов буквы ё. Верстальщик за этим не следил)
\item Уделить внимание самим аккордам. Не факт что все песни
в правильных/удобных тональнастях, нужно призвать для
этого гитаристов
\item общий осмотр: поиск ошибок содержания, логики, эстетики
\end{enumerate}
\newpage
Версия песенника 3.3
\begin{enumerate}
\item Версия включает исправления по списку песен 28.03.2023
\item Версия включает исправления редакции от 27.04.2023. \\ docs.google.com/document/d/1ccpFXUHAUCM9uB7PMvuHNoh-e1bhCoorUNitSKM8Mm4. 
\item Весь текст песенника был проверен на ресурсе online.orfo.ru.
\item исправлены аккорды песен: Луч солнца золотого
\item Перестроен порядок жаворонков: Непогода и Мамонты поменялись местами; Большой секрет сдвинут через 4 песни
\item Унифицировались знаки пунктуации:
\begin{itemize}
    \item Все угловатые кавачки «» заменены на `` '' кроме последней страницы -- они там к месту
    \item Все знаки тире которые обозначались ``-'' стали ``--''.
    \item Все вхождения ``. . .'' были заменены на ``...'' (баг latex с символом троеточия)
    \item Все символы ``|'' убраны в проигрышах
\end{itemize} 
\item Исправлены оступы в песнях где есть указание  ``2 раза'' через фигурные скобки (Этот баг не исправлен! я вручную передвинул каждый текст на отрицательное значение вниз)
\item Улучшен вид табулатур (пришлось переписать библиотеку)
\end{enumerate}
Необходимо далее:
\begin{enumerate}
\item Уделить внимание самим аккордам. Не факт что все песни в правильных/удобных тональнастях, нужно призвать для этого гитаристов
\item общий осмотр: поиск ошибок содержания, логики, эстетики
\end{enumerate}
\newpage
Версия песенника 3.4
\begin{enumerate}
\item Версия включает исправления по списку песен 23.04.2024
\item Переставлены некоторые песни для сохранения чётности и отображения длинных песен на развороте.
\item откорректированны ссылки на страницы песен в самоучителе (стр 100-101) согласно правкам выше
\end{enumerate}
Замечания верстальщика: 
\begin{enumerate}
\item песня Рыжий кошка сашка абсолютно изменены аккорды, отличаются от тех которые были даны
\item песня Меняй Павел Пиковский взят аккорд Си-минор, что (вроде) отличается от тех аккорда Си который дан
\item все вхождения B аккордов заменены на H, так принято в этом песеннике
\end{enumerate}
Необходимо далее:
\begin{enumerate}
\item Уделить внимание аккордам. Не факт что все песни в правильных/удобных тональнастях, нужно призвать для этого гитаристов
\item общий осмотр: поиск ошибок содержания, логики, эстетики
\end{enumerate}
\newpage
\end{document}